% Einheit 2 – Kreisfläche – Nr. 22–32
\setcounter{aufgabe}{21}
\einheitenkopf[e2][2 Kreisfläche]{Einheit 2 von 3 · Kreisfläche}
\zweigzeile{Hier lernst du, die Fläche eines Kreises zu berechnen und aus der Fläche den Radius · neu oder schon bekannt – je nach Buch · P10 oft · baut auf: Kreisumfang (Einheit 1), Quadrieren und Wurzelziehen, Formeln umstellen}

\begin{aufgabe}{Ich erkenne, welche Formel passt.}
Kreuze an, welche Formel du brauchst. Rechne nicht.
\begin{teile}
\teil Wie viel Stoff braucht man für eine runde Tischdecke? \hfill \kreuz{$u = \pi \cdot d$}\kreuz{$A = \pi \cdot r^2$}
\teil Wie lang ist die Kante eines runden Tellers? \hfill \kreuz{$u = \pi \cdot d$}\kreuz{$A = \pi \cdot r^2$}
\teil Wie viel Platz nimmt ein runder Teppich ein? \hfill \kreuz{$u = \pi \cdot d$}\kreuz{$A = \pi \cdot r^2$}
\teil Wie lang muss eine Lichterkette einmal um einen Baumstamm sein? \hfill \kreuz{$u = \pi \cdot d$}\kreuz{$A = \pi \cdot r^2$}
\teil Wie viel Farbe braucht man für ein rundes Schild? \hfill \kreuz{$u = \pi \cdot d$}\kreuz{$A = \pi \cdot r^2$}
\end{teile}
\end{aufgabe}

\verfahren{Fläche aus Radius oder Durchmesser}
\begin{aufgabe}{Ich kann die Fläche eines Kreises berechnen.}
Berechne die Fläche $A$ aus dem Radius $r$. Runde auf eine Stelle nach dem Komma.
\begin{teilezwei}
\tz $r = 1\,\text{cm}$ \quad \feld[cm²]{A} & \tz $r = 6\,\text{cm}$ \quad \feld[cm²]{A} \\
\tz $r = 7\,\text{cm}$ \quad \feld[cm²]{A} & \tz $r = 10\,\text{m}$ \quad \feld[m²]{A} \\
\anweisung{Jetzt ist der Durchmesser $d$ gegeben.}
\tz $d = 18\,\text{cm}$ \quad \feld[cm²]{A} & \tz $d = 5{,}4\,\text{cm}$ \quad \feld[cm²]{A} \\
\end{teilezwei}
\end{aufgabe}

\begin{aufgabe}{Ich kann die Fläche von Halbkreis und Viertelkreis berechnen.}
Berechne die Fläche. Runde auf eine Stelle nach dem Komma.
\begin{teile}
\teil Halbkreis mit $r = 6\,\text{cm}$ \quad \leerfeld[cm²]
\teil Viertelkreis mit $r = 12\,\text{cm}$ \quad \leerfeld[cm²]
\teil Halbkreis mit $d = 14\,\text{cm}$ \quad \leerfeld[cm²]
\teil Dreiviertelkreis mit $r = 4\,\text{cm}$ \quad \leerfeld[cm²]
\end{teile}
\end{aufgabe}

\verfahren{Flächenterme zu Kreisfiguren}
\begin{aufgabe}{Ich kann zu einer Kreisfigur den Term für die Fläche aufschreiben.}
Schreibe den Term für die graue Fläche auf. Rechne nicht aus.

\begin{kreis}[1.4]\mittelpunkt{}\sektor{0}{90}{}\radius{0}{3\,\text{cm}}\end{kreis}\hfill
\begin{kreis}[1.4]\mittelpunkt{}\sektor{0}{180}{}\durchmesser{0}{8\,\text{cm}}\end{kreis}\hfill
\begin{kreis}[1.4]\mittelpunkt{}\sektor{0}{270}{}\radius{270}{2\,\text{cm}}\end{kreis}

\begin{teile}
\teil linke Figur: \quad $A =$ \leerfeld
\teil mittlere Figur: \quad $A =$ \leerfeld
\teil rechte Figur: \quad $A =$ \leerfeld
\end{teile}
\end{aufgabe}

\begin{aufgabe}{Ich kann zu einem Flächenterm die Kreisfigur angeben.}
Gib an, welcher Teil vom Kreis mit welchem Radius zum Term gehört.
\begin{teile}
\teil $\frac{1}{2} \cdot \pi \cdot 5^2$ \hfill \leerfeld \quad \leerfeld
\teil $\frac{1}{4} \cdot \pi \cdot 6^2$ \hfill \leerfeld \quad \leerfeld
\end{teile}
\end{aufgabe}

\verfahren{Radius, Durchmesser, Umfang und Fläche}
\begin{aufgabe}{Ich kann in einer Tabelle Radius, Durchmesser, Umfang und Fläche ergänzen.}
Ergänze die Tabelle. Runde auf eine Stelle nach dem Komma.

\sachtabelle{lcccc}{ & Radius $r$ & Durchmesser $d$ & Umfang $u$ & Fläche $A$}{Kreis 1 & $6\,\text{m}$ & \leerzelle & \leerzelle & \leerzelle \\ Kreis 2 & \leerzelle & $9\,\text{cm}$ & \leerzelle & \leerzelle \\ Kreis 3 & \leerzelle & $2{,}4\,\text{m}$ & \leerzelle & \leerzelle \\ Kreis 4 & \leerzelle & \leerzelle & $22{,}0\,\text{cm}$ & \leerzelle}

\end{aufgabe}

\begin{aufgabe}{Ich kann zwischen Fläche, Radius und Durchmesser umrechnen.}
Berechne den Radius $r$ aus der Fläche $A$. Runde auf eine Stelle nach dem Komma.
\begin{teilezwei}
\tz $A = 154\,\text{cm}^2$ \quad \feld[cm]{r} & \tz $A = 314\,\text{m}^2$ \quad \feld[m]{r} \\
\anweisung{Berechne den Durchmesser $d$.}
\tz $A = 20\,\text{cm}^2$ \quad \feld[cm]{d} & \\
\end{teilezwei}
\anweisung{Runde auf zwei Stellen nach dem Komma.}
\begin{teile}
\teil Ein runder Brunnen hat eine Wasserfläche mit dem Durchmesser $3{,}46\,\text{m}$. Berechne den Inhalt der Wasserfläche. (P10 2016 OS) \hfill \leerfeld[m²]
\end{teile}
\end{aufgabe}

\begin{aufgabe}{Ich finde den Fehler in einer Flächenrechnung.}
Finde den Fehler, benenne ihn und korrigiere die Rechnung.
\begin{teile}
\teil Kreis mit dem Durchmesser $d = 12\,\text{cm}$: \rechnung{A &= \pi \cdot 12^2 \\ &\approx 452{,}4\,\text{cm}^2}
\teil Kreis mit dem Radius $r = 9\,\text{m}$: \rechnung{A &= 2 \cdot \pi \cdot 9 \\ &\approx 56{,}5\,\text{m}^2}
\teil Kreis mit dem Radius $r = 7\,\text{cm}$: \rechnung{A &= \pi \cdot 7^2 \\ &= \pi \cdot 14 \\ &\approx 44{,}0\,\text{cm}^2}
\end{teile}
\end{aufgabe}

\begin{aufgabe}{Ich kann begründen, wie die Fläche vom Radius abhängt.}
Begründe.
\begin{teile}
\teil Der Radius eines Kreises wird verdoppelt. Mia sagt: „Dann verdoppelt sich auch die Fläche.“ Hat sie recht?
\teil Ein Kreis wird in viele schmale Stücke geschnitten. Die Stücke werden abwechselnd mit der Spitze nach oben und nach unten nebeneinandergelegt, so entsteht fast ein Rechteck. Warum ist es ungefähr so lang wie der halbe Umfang und so hoch wie der Radius? Warum folgt daraus $A = \pi \cdot r^2$?
\end{teile}
\end{aufgabe}

\begin{aufgabe}{Ich kann unterscheiden, ob Umfang oder Fläche gesucht ist.}
Entscheide, was gesucht ist, und berechne es. Runde auf eine Stelle nach dem Komma.
\begin{teile}
\teil Ein runder Tisch hat den Durchmesser $1{,}2\,\text{m}$. Wie viel Stoff braucht man für eine Tischdecke, die genau die Tischplatte bedeckt? \hfill \leerfeld
\teil Wie lang muss eine Borte sein, die einmal um den Rand dieses Tisches reicht? \hfill \leerfeld
\teil Ein rundes Beet hat den Radius $3{,}5\,\text{m}$. Wie viele Meter Kantensteine braucht man für den Rand? \hfill \leerfeld
\teil Wie groß ist die Fläche, die im Beet mit Rasen eingesät wird? \hfill \leerfeld
\end{teile}
\end{aufgabe}

\begin{aufgabe}{Ich kann mit Kreisflächen Preise vergleichen.}
Eine große Pizza mit $30\,\text{cm}$ Durchmesser kostet $8\,€$. Zwei kleine Pizzen mit je $20\,\text{cm}$ Durchmesser kosten zusammen $9\,€$.
\begin{teile}
\teil Berechne die Fläche der großen Pizza und die Fläche der beiden kleinen zusammen. Runde auf eine Stelle.
\teil Berechne für beide Angebote den Preis für $100\,\text{cm}^2$ Pizza. Runde auf Cent.
\teil Welches Angebot ist günstiger? Begründe mit deinen Ergebnissen.
\end{teile}
\end{aufgabe}
